\section{Analytical Model and Numerical Method}

\subsection{Disk model}

We consider an embryo $i$ interior to a single giant planet $j$. Their
semimajor axes, eccentricities, masses, and mean motions are denoted by
$(a_i,e_i,M_i,n_i)$ and $(a_j,e_j,M_j,n_j)$, respectively, and the stellar
mass is $M_*$. The axisymmetric gas disk has surface density
\begin{equation}
    \Sigma(r,t)=\Sigma_1
    \left(\frac{r}{1\,\mathrm{AU}}\right)^{-k}u(t),
    \qquad
    u(t)=\exp\left(-\frac{t}{T_{\rm dep}}\right).
\label{eq:sigma_analytical}
\end{equation}
Because disk gravity is linear in $\Sigma$, every disk-induced precession
rate can be written as $g_d(a,t)=u(t)g_{d,0}(a)$.

For an untruncated power-law disk, the radial acceleration is
\begin{equation}
    F_{d,0}(r,t)=-\frac{2\pi G\Sigma(r,t)}{2-k}Z_k,
\label{eq:full_disk_force}
\end{equation}
where
\begin{eqnarray}
    Z_k &=& 1+(2-k)(k-1)
    \sum_{\ell=1}^{\infty}
    \frac{(4\ell+1)A_\ell}
    {(2\ell-1+k)(2\ell+2-k)},
    \nonumber\\
    A_\ell&=&
    \left[\frac{(2\ell)!}{2^{2\ell}(\ell!)^2}\right]^2.
\label{eq:zk_coefficients}
\end{eqnarray}
The corresponding apsidal precession of a nearly circular orbit is
\begin{equation}
    g_{d,0}(a)=-\frac{\pi G\Sigma(a,0)}{n(a)a}Z_k,
    \qquad
    n(a)=\left(\frac{GM_*}{a^3}\right)^{1/2}.
\label{eq:full_disk_precession}
\end{equation}
For the default $k=3/2$ disk, $Z_k=1.0942$.

To describe a truncated disk, we define the acceleration from the disk
exterior to an edge $R$ for $r<R$,
\begin{equation}
    F_{>R}(r,t)=2\pi G\Sigma(r,t)
    \sum_{\ell=1}^{\infty}
    \frac{2\ell A_\ell}{2\ell-1+k}
    \left(\frac{r}{R}\right)^{2\ell-1+k},
\label{eq:exterior_disk_force}
\end{equation}
and the acceleration from the disk interior to $R$ for $r>R$,
\begin{eqnarray}
    F_{<R}(r,t)&=&-2\pi G\Sigma(r,t)
    \left\{\frac{(R/r)^{2-k}}{2-k}\right.\nonumber\\
    &&\left.+\sum_{\ell=1}^{\infty}
    \frac{(2\ell+1)A_\ell}{2\ell+2-k}
    \left(\frac{R}{r}\right)^{2\ell+2-k}\right\}.
\label{eq:interior_disk_force}
\end{eqnarray}
The first term of Equation~\ref{eq:interior_disk_force} is the
$\ell=0$ monopole. It contributes to the radial force but not to apsidal
precession. For any small axisymmetric radial perturbation $F(r)$, the
precession rate used here is
\begin{equation}
    g_F(a)=\frac{1}{2n(a)}
    \left[\frac{2F}{a}+\frac{dF}{da}\right].
\label{eq:force_to_precession}
\end{equation}
It is therefore useful to define
\begin{eqnarray}
    {\cal E}(a;R)&=&
    \sum_{\ell=1}^{\infty}
    \frac{2\ell(2\ell+1)A_\ell}{2\ell-1+k}
    \left(\frac{a}{R}\right)^{2\ell-1+k},
    \nonumber\\
    {\cal I}(a;R)&=&
    \sum_{\ell=1}^{\infty}
    \frac{2\ell(2\ell+1)A_\ell}{2\ell+2-k}
    \left(\frac{R}{a}\right)^{2\ell+2-k},
\label{eq:edge_sums}
\end{eqnarray}
and $P(a)=\pi G\Sigma(a,0)/[n(a)a]$. Thus, an exterior disk gives
$g_{>R}=P{\cal E}$, whereas an interior disk gives $g_{<R}=P{\cal I}$.

Let $h_J=(M_j/3M_*)^{1/3}$. Following the reference implementation, the
one-Hill-radius gap edges are
\begin{equation}
    (R_{\rm in},R_{\rm out})=
    \begin{cases}
    \left(a_j(1-e_j)(1-h_J),a_j(1+e_j)(1+h_J)\right),&e_j>h_J,\\
    \left(a_j(1-h_J),a_j(1+h_J)\right),&e_j\leq h_J.
    \end{cases}
\label{eq:hill_gap_edges}
\end{equation}
The edges are fixed from the assigned initial giant orbit. Inner embryos use
the $R_{\rm in}$ correction from the Multi\_HR5183 implementation, while
bodies inside the gap use the two truncated disks,
\begin{equation}
g_{\rm empty}(a)=
\begin{cases}
-2P\left[(2-k)Z_k+\displaystyle\sum_{\ell=1}^{9}
\frac{\ell(2\ell+1)A_\ell}{2\ell-1+k}
\left(\frac{a}{R_{\rm in}}\right)^{2\ell-1+k}\right],
&a<R_{\rm in},\\[3pt]
 P\left[{\cal I}(a;R_{\rm in})
             +{\cal E}(a;R_{\rm out})\right],
&R_{\rm in}<a<R_{\rm out},\\[3pt]
-PZ_k,&a\geq R_{\rm out}.
\end{cases}
\label{eq:empty_gap_precession}
\end{equation}
If the gap retains a fraction $k_{\rm gap}$ of the unperturbed density,
linearity gives
\begin{equation}
    g_{\rm depgap}(a)=g_{\rm empty}(a)
    +k_{\rm gap}\left[g_{d,0}(a)-g_{\rm empty}(a)\right].
\label{eq:depleted_gap_precession}
\end{equation}
Here $k_{\rm gap}=0$ is an empty gap and $k_{\rm gap}=1$ recovers the
no-gap disk. The default is the empty-gap limit, $k_{\rm gap}=0$. The same
interpolation is used for the radial acceleration.
An optional nonzero disk inner edge is implemented by subtracting the
contribution of the disk interior to that edge in the general truncated-disk
branches; the restored legacy $a<R_{\rm in}$ branch does not include this
additional subtraction. The default model has no inner cutoff.

\subsection{Location and time of the secular resonance}

The giant causes the embryo to precess at the Laplace--Lagrange rate
\begin{equation}
    g_{i,j}=\frac{G M_j\alpha b_{3/2}^{(1)}(\alpha)}
    {4n_i a_i^2a_>},
    \qquad
    \alpha=\frac{a_<}{a_>},
\label{eq:gij_analytical}
\end{equation}
where $a_<\equiv\min(a_i,a_j)$ and $a_>\equiv\max(a_i,a_j)$. We evaluate the
Laplace coefficient numerically,
\begin{equation}
    b_s^{(m)}(\alpha)=\frac{1}{\pi}
    \int_0^{2\pi}
    \frac{\cos(m\phi)\,d\phi}
    {(1+\alpha^2-2\alpha\cos\phi)^s},
\label{eq:laplace_integral}
\end{equation}
rather than using a truncated power series. General relativity adds the
low-eccentricity precession
\begin{equation}
    g_{i,\rm GR}=\frac{3n_iGM_*}{a_i c^2}.
\label{eq:gr_precession}
\end{equation}

The secular resonance occurs where the embryo and giant have equal
apsidal precession rates,
\begin{equation}
    g_{i,j}+g_{i,\rm GR}+u(t_R)g_{i,d,0}
    =u(t_R)g_{j,d,0}.
\label{eq:sr_condition}
\end{equation}
The disk precession of each body is evaluated from the same disk geometry;
in particular, the giant lies between the two edges of its own gap and is
assigned the middle branch of Equation~\ref{eq:empty_gap_precession} or
its partially depleted counterpart. The small GR precession of the giant
is neglected in the analytical resonance location. Solving
Equation~\ref{eq:sr_condition} gives
\begin{equation}
    u_R(a_i)=\exp\left[-\frac{t_R(a_i)}{T_{\rm dep}}\right]
    =\frac{g_{i,j}+g_{i,\rm GR}}
    {g_{j,d,0}-g_{i,d,0}},
\label{eq:u_r_single_giant}
\end{equation}
and
\begin{equation}
    \frac{t_R(a_i)}{T_{\rm dep}}=-\ln u_R(a_i).
\label{eq:t_r_single_giant}
\end{equation}
Only $0<u_R\leq1$ corresponds to a resonance reached during disk
depletion. Equations~\ref{eq:u_r_single_giant} and
\ref{eq:t_r_single_giant} therefore specify both where and when the
secular resonance sweeps through the inner system.

\subsection{Estimate of the surfing regime}

The scale height is prescribed as
\begin{equation}
    H(a)=H_0\left(\frac{a}{1\,\mathrm{AU}}\right)^\beta,
\end{equation}
and the eccentricity-damping time of an embryo is
\begin{equation}
    t_e(a,t)=\frac{1}{f_e}
    \left(\frac{M_*}{M_i}\right)
    \left(\frac{M_*}{\Sigma(a,t)a^2}\right)
    \left(\frac{H}{a}\right)^4\Omega_K^{-1},
    \quad
    \Omega_K=\left(\frac{GM_*}{a^3}\right)^{1/2}.
\label{eq:type1_damping_time}
\end{equation}
At resonance passage we use $\Sigma(a,t_R)=\Sigma(a,0)u_R(a)$. We prescribe
angular-momentum-conserving eccentricity damping with the radial force
\begin{equation}
    \mathbf a_e=-\frac{2v_r}{t_e(a,t_R)}\hat{\mathbf r}.
\label{eq:radial_eccentricity_damping_force}
\end{equation}
The factor of two makes $t_e$ the eccentricity e-folding time in the
small-eccentricity limit. Because the force is radial,
\begin{equation}
    \dot{\mathcal E}=-\frac{2v_r^2}{t_e},
    \qquad \dot h=0.
\label{eq:damping_energy_angular_momentum}
\end{equation}
On a Kepler orbit,
$\langle v_r^2\rangle=(GM_*/a)[1-\sqrt{1-e^2}]$. Therefore
\begin{equation}
    \frac{\left\langle\dot a_{\rm circ}\right\rangle}{a}
    =-\frac{4}{t_e(a,t_R)}\left(1-\sqrt{1-e^2}\right),
\label{eq:circularization_speed}
\end{equation}
with the small-eccentricity limits
\begin{equation}
    \frac{\langle\dot e\rangle}{e}=-\frac{1}{t_e}+O(e^2),
    \qquad
    \frac{\langle\dot a_{\rm circ}\rangle}{a}
    =-\frac{2e^2}{t_e}+O(e^4).
\label{eq:radial_damping_orbit_average}
\end{equation}

The resonance trajectory has slope
\begin{equation}
    \frac{dt_R}{da}=-T_{\rm dep}\frac{d\ln u_R}{da}.
\label{eq:sr_slope}
\end{equation}
Because the damping-driven migration is inward, compatible surfing requires
$dt_R/da<0$, or equivalently $d\ln u_R/da>0$. For an eccentricity
$e_{\rm kick}$ generated at resonance, the surfing map evaluates the
dimensionless speed ratio
\begin{eqnarray}
    {\cal S}(a,T_{\rm dep})
    &=&|\dot a_{\rm circ}|
       \left|\frac{dt_R}{da}\right|
    \nonumber\\
    &=&\frac{4a[1-\sqrt{1-e_{\rm kick}^2}]}
    {t_e(a,t_R)}
    T_{\rm dep}\left|\frac{d\ln u_R}{da}\right|.
\label{eq:surfing_number}
\end{eqnarray}
Thus ${\cal S}=1$ means that circularization-driven migration and the resonance sweep
have equal radial speeds. The corresponding depletion-time boundary is
\begin{equation}
    T_{\rm dep,S=1}(a)=
    \frac{t_e(a,t_R)}
    {4a[1-\sqrt{1-e_{\rm kick}^2}]
    |d\ln u_R/da|}.
\label{eq:surfing_boundary}
\end{equation}
Values ${\cal S}\gtrsim1$ identify inward-sweeping configurations in which
migration can be fast enough to follow the resonance. Sustained capture is
not guaranteed by this local speed comparison alone.

We also require resonance passage to be capable of producing the assumed
$e_{\rm kick}$. The eccentricity forcing coefficient is
\begin{equation}
    B_{i,j}=\frac{n_i}{4}\frac{M_j}{M_*}
    \alpha^2b_{3/2}^{(2)}(\alpha),
\label{eq:secular_forcing_coefficient}
\end{equation}
and the stationary-phase estimate for a linear resonance crossing is
\begin{equation}
    e_{\rm sp}\simeq |B_{i,j}|e_j
    \left(\frac{2\pi}{|\dot{\Delta g}_R|}\right)^{1/2},
\label{eq:stationary_phase_kick}
\end{equation}
where
\begin{equation}
    \Delta g=g_i-g_j,
    \qquad
    |\dot{\Delta g}_R|
    =\frac{u_R}{T_{\rm dep}}
    |g_{j,d,0}-g_{i,d,0}|.
\label{eq:relative_precession_sweep}
\end{equation}
The analytical surfing candidates satisfy both ${\cal S}\gtrsim1$ and
$e_{\rm sp}\gtrsim e_{\rm kick}$. The second condition is an undamped
linear estimate, so the resulting map is a configuration search rather
than a proof of resonant capture.

The current configuration search adopts the NLT05 parameters: a
$1\,M_{\rm Jup}$ giant at $5.203\,\mathrm{AU}$,
$\Sigma_1=1700\,\mathrm{g\,cm^{-2}}$, $k=3/2$,
$H_0=0.05\,\mathrm{AU}$, $\beta=1.25$, and $f_e=1$. It evaluates embryo
masses $0.03$, $0.3$, $1$, and $3\,M_\oplus$ over
$0.4\leq a_i/\mathrm{AU}\leq3$ and
$10^5\leq T_{\rm dep}/\mathrm{yr}\leq3\times10^7$, and shows both the
no-gap and empty-gap limits. GR is disabled for this particular parameter
survey. These choices define the plotted search grid, rather than a
restriction of Equations~\ref{eq:surfing_number}--
\ref{eq:relative_precession_sweep}.

\subsection{Numerical setup and prescribed forces}

The $N$-body calculations are performed with REBOUND in units of yr, AU,
and $M_\odot$, for which $G=4\pi^2$. Newtonian gravity is integrated with
MERCURIUS. The additional acceleration on a nonstellar body is
\begin{equation}
    \mathbf a_{\rm add}=F_d(r,t)\hat{\mathbf r}
    +\mathbf a_{\rm GR}+\mathbf a_e+\mathbf a_a.
\label{eq:additional_acceleration}
\end{equation}
The disk acceleration is applied to both the giant and the embryos. With the
Multi\_HR5183 inner-embryo correction restored, it is
\begin{equation}
F_{\rm empty}=\begin{cases}
-4\pi G\Sigma(r,t)\left[z_k+\displaystyle\sum_{\ell=1}^{9}
\frac{\ell A_\ell}{2\ell-1+k}
\left(\frac{r}{R_{\rm in}}\right)^{2\ell-1+k}\right],&a<R_{\rm in},\\
F_{<R_{\rm in}}+F_{>R_{\rm out}},&R_{\rm in}<a<R_{\rm out},\\
F_{d,0},&a\geq R_{\rm out}.
\end{cases}
\label{eq:empty_gap_force}
\end{equation}
The partially depleted force is extended by linear interpolation,
\begin{equation}
    F_{\rm depgap}=F_{\rm empty}
    +k_{\rm gap}(F_{\rm full}-F_{\rm empty}).
\label{eq:depleted_gap_force}
\end{equation}
The gap edges are calculated once from the assigned initial $a_j$ and $e_j$
using Equation~\ref{eq:hill_gap_edges}. The general edge series are truncated
at $\ell=30$ in the default calculation, while the restored Multi\_HR5183
inner-embryo correction retains its original $\ell=1,\ldots,9$ truncation.

The first post-Newtonian stellar acceleration is applied to embryos but not
to the gap-opening giant,
\begin{equation}
    \mathbf a_{\rm GR}
    =\frac{GM_*}{c^2r^3}
    \left[\left(\frac{4GM_*}{r}-v^2\right)\mathbf r
    +4(\mathbf r\mathbin{\cdot}\mathbf v)\mathbf v\right],
\label{eq:numerical_gr_force}
\end{equation}
with $c=63241.077\,\mathrm{AU\,yr^{-1}}$.

The dissipative forces are applied only to embryos. At each force
evaluation, $\Sigma$, $H$, $\Omega_K$, and the remaining factors in the
damping time are evaluated at the osculating semimajor axis $a$. Defining
\begin{equation}
    t_{\rm damp}=\left(\frac{M_*}{M_i}\right)
    \left(\frac{M_*}{\Sigma(a,t)a^2}\right)
    \left(\frac{H(a)}{a}\right)^4\Omega_K^{-1},
\label{eq:numerical_damping_time}
\end{equation}
which is Equation~\ref{eq:type1_damping_time} before applying $f_e$, the
tidal damping force is
\begin{equation}
    \mathbf a_e=-\frac{2[(\mathbf v_i-\mathbf v_*)\mathbin{\cdot}
    \hat{\mathbf r}]}{t_e}\hat{\mathbf r},
    \qquad t_e=\frac{t_{\rm damp}}{f_e}.
\label{eq:numerical_eccentricity_force}
\end{equation}
This radial force exerts no torque, so eccentricity damping and explicit
Type-I migration remain independently selectable.
The default $f_e=1$ gives the nominal analytical damping strength, and
$f_e=0$ disables this force.

An optional explicit Type-I semimajor-axis force can also be enabled. We
use the Tanaka et al. locally isothermal migration time for
$\Sigma\propto r^{-k}$,
\begin{equation}
    t_a=\frac{1}{f_a(2.7+1.1k)}
    \left(\frac{M_*}{M_i}\right)
    \left(\frac{M_*}{\Sigma a^2}\right)
    \left(\frac{H}{a}\right)^2\Omega_K^{-1},
\label{eq:numerical_migration_time}
\end{equation}
and prescribe the associated inward tangential acceleration as
\begin{equation}
    \mathbf a_a=-\frac{v_K(r)}{2t_a}
    \frac{\mathbf v-(\mathbf v\mathbin{\cdot}\hat{\mathbf r})
    \hat{\mathbf r}}
    {|\mathbf v-(\mathbf v\mathbin{\cdot}\hat{\mathbf r})
    \hat{\mathbf r}|}.
\label{eq:numerical_migration_force}
\end{equation}
The baseline $\mathtt{open\_typeI}=1$ uses only the radial eccentricity-damping
force. Setting $\mathtt{open\_typeI}=2$ additionally enables the explicit
Type-I torque. The surfing estimate above uses only
Equation~\ref{eq:circularization_speed}; it does not include this explicit
torque.

The default disk parameters in the numerical force are
$\Sigma_1=1700\,\mathrm{g\,cm^{-2}}$, $k=3/2$,
$H_0=0.025\,\mathrm{AU}$, $\beta=1.25$,
$T_{\rm dep}=1\,\mathrm{Myr}$, $N_H=1$, and $k_{\rm gap}=0$.
The timestep and output interval are configurable. The baseline source
sets the MERCURIUS timestep to $0.1$ times the innermost embryo's initial
period; run scripts can override this value. Collisions are detected
using physical radii and resolved as perfect mergers, and the full
simulation state is stored in a REBOUND Simulationarchive for subsequent
orbital and precession analysis.
